\documentclass[11pt,a4paper]{article}
\usepackage[margin=26mm]{geometry}
\usepackage{amsmath,amssymb,amsthm,hyperref}
\newtheorem{proposition}{Proposition}
\title{Checking intersection properties on intervals}
\author{Human Compression project}\date{}
\begin{document}\maketitle
This note records an exact implementation of the end-intersection criterion
already used in the acquired permutation enumeration~\cite{method}.
The new verifier starts with a complete literal tuple. It imports no
prefix certificate from the enumeration and does not alter that search.

Let $t=(t_1,\ldots,t_r)$ be a tuple of permutations, with $r\geq1$.
For $1\leq a\leq b\leq r$, write
$H_{a,b}=\langle t_a,\ldots,t_b\rangle$.
The intersection property means that
\[
 \langle t_i:i\in I\rangle\cap\langle t_j:j\in J\rangle
 =\langle t_k:k\in I\cap J\rangle
 \quad(I,J\subseteq\{1,\ldots,r\}).
\]

\begin{proposition}
The tuple $t$ is a string C-marking if and only if its entries are distinct
nonidentity involutions, its distant entries commute, and
\begin{equation}\label{eq:ends}
 H_{a,b-1}\cap H_{a+1,b}=H_{a+1,b-1}
 \qquad(1\leq a<b-1\leq r-1).
\end{equation}
For an accepted tuple, the last test uses exactly
$(r-1)(r-2)/2$ actual subgroup intersections.
\end{proposition}
\begin{proof}
Necessity follows by applying the full intersection property to the two
displayed intervals. Distinctness follows by applying it to distinct
singleton index sets. The string relations and involution requirements
are part of the definition.

For sufficiency, induct on the length of a contiguous interval.
A single nonidentity involution gives the rank-one marking.
Two distinct involutions give the rank-two marking: their order-two
subgroups have trivial intersection, which is the only nontrivial
intersection condition in that rank. This includes the case where the
two involutions commute.
For an interval of length at least three, both end omissions are string
C-markings by the induction hypothesis. Their actual intersection is
their common middle subgroup by~\eqref{eq:ends}.
The acquired end-intersection criterion~\cite[Section~1]{method} now gives
the intersection property on that whole interval. Its string relations
are restrictions of the originally checked relations. This proves the
induction and hence the result for $[1,r]$.
There are $r-\ell+1$ intervals of length $\ell$, so the number of tests is
$\sum_{\ell=3}^r(r-\ell+1)=(r-1)(r-2)/2$.
\end{proof}

The routine constructs the actual subgroups $H_{a,b}$ in increasing
interval length. All three groups in each comparison have therefore
already been constructed. It tests equality of subgroups, rather than
equality of their orders or the intersection of projected images.
The full generating-group order is a separate check in an application;
the proposition applies equally to proper subgroups and to disconnected
string diagrams. It makes no classification or counting assertion.

\section*{Controls and scope}
The independent comparison routine constructs every subgroup associated
with an arbitrary index subset. It compares the two decisions on all
ordered tuples of one through four involutions of $S_4$, including tuples
with repeated entries and tuples failing string commutation. No tuple
is removed from that comparison by the new criterion. Five explicit
invalid inputs also cover the empty tuple, the identity, an element of
order three, a repeated mark and a distant noncommuting pair.

The retained complete $S_8$, $S_{10}$ and $S_{11}$ positive lists are
read directly from their accepted full-subset certificates. The new
routine checks every tuple, including its actual symmetric generating
group. The extraction is compared literally against those frozen
certificates; the underlying enumeration is not repeated. These controls
test the implementation. The proposition, with its acquired criterion,
justifies the replacement of exhaustive subset checks in later uses.
Permutation arithmetic and exact subgroup intersections in GAP remain
explicit computational infrastructure.

The end-intersection criterion is borrowed, and its original provenance
is retained through~\cite{method}. This component only reorganizes its
evaluation and supplies a separate verifier. It establishes no additional
survey cell, whole-original recovery, open conjecture or compression rate.
Any application must separately pay completeness, target identification,
the source's isomorphism convention and the exact numerical output.

\begin{thebibliography}{9}\small
\bibitem{method} Human Compression project, \emph{Complete permutation
searches and the six-point row}, accepted finite-method component,
Section~1, with its independent review and root adoption.
\end{thebibliography}
\end{document}
